\chapter{Critical Analysis of the Four Source Papers}
\label{ch:sources}

The project began from four papers supplied with the topic. This chapter
reads each in turn and then asks what they show together. The analysis used
the full text and reference lists. Because text extraction lost most of
their tables, every table was read again from the rendered PDF pages. Where
a paper's text and tables disagree, both values are reported and the
disagreement is listed in Table~\ref{tab:inconsistencies}.

One identification matters from the start. The fourth file, supplied as
\emph{Cross-Modality Fusion Mamba for Object Detection}, is CFMW by Li and
colleagues~\citep{li2025cfmw}. It is not the similarly named Fusion-Mamba by
Dong and colleagues~\citep{dong2025fusionmamba}. Both are cited throughout,
and both belong to the same family of designs. Table~\ref{tab:srcoverview}
summarises the four papers.

\begin{table}[htbp]
\centering
\caption[The four source papers at a glance]{The four source papers at a glance.}
\label{tab:srcoverview}
\small
\begin{tabularx}{\textwidth}{@{}P{2.25cm}P{2.75cm}P{2.55cm}LP{1.45cm}@{}}
\hd{Paper} & \hd{Venue} & \hd{Base detector} & \hd{Task and data} & \hd{Code} \\ \gap
MambaVision \citep{hatamizadeh2025mambavision} & CVPR 2025 (NVIDIA) & Cascade Mask R-CNN & backbone; ImageNet, COCO, ADE20K & public \\ \sgap
HPS-DETR \citep{wang2025hpsdetr} & IEEE TGRS 63, 2025 & RT-DETR-r18 & aerial small objects; SIMD, DOTA, NWPU VHR-10 & none \\ \sgap
UAVDet \citep{yang2026uavdet} & CVIU 264, 2026 & RTMDet-m & UAV small objects; VisDrone, UAVDT, DroneVehicle (RGB) & none \\ \sgap
CFMW \citep{li2025cfmw} & IEEE TCSVT 35(12), 2025 & dual-stream YOLOv5-L & visible--thermal under weather; LLVIP, SWVI & announced \\
\end{tabularx}
\end{table}

\section{MambaVision: a hybrid Mamba--Transformer backbone}
\label{sec:src-mv}

\paragraph{What it proposes}
MambaVision, by Hatamizadeh and Kautz of NVIDIA, is a four-stage
hierarchical backbone~\citep{hatamizadeh2025mambavision}. Stages one and two,
at strides 4 and 8, are plain convolutional residual blocks chosen for speed
at high resolution. Stages three and four, at strides 16 and 32, contain $N$
layers each. The first $N/2$ use a redesigned Mamba mixer and the last $N/2$
use windowed self-attention. The mixer departs from Mamba in three ways.
\begin{itemize}
  \item The causal convolution is replaced by an ordinary one.
  \item A second branch is added that does not pass through the SSM.
  \item The multiplicative gate is replaced by concatenating the two
        branches, followed by a linear projection.
\end{itemize}
The scan runs once, in one direction.

\paragraph{Evidence}
Table~\ref{tab:mvcoco} gives the ImageNet and COCO results that matter here.
The ablations support both design choices. Concatenation instead of gating
adds 1.0 point of ImageNet top-1 and 1.1 box AP with Mask R-CNN. Placing
attention in the last half of each stage gives 82.3\% top-1. That beats
random placement (81.3\%), the first half (81.5\%) and alternating patterns
(81.4--81.6\%).

\begin{table}[htbp]
\centering
\caption[MambaVision: classification throughput and COCO detection]{MambaVision
results reproduced from the paper. Throughput is images per second at
$224\times224$, batch 128, on an A100. COCO uses Cascade Mask R-CNN with the
$3\times$ schedule. Parameters and GFLOPs for COCO refer to the whole
detector.}
\label{tab:mvcoco}
\small
\begin{tabularx}{\textwidth}{@{}LR{1.55cm}R{1.35cm}R{1.6cm}R{1.6cm}R{1.6cm}@{}}
\hd{Model} & \hd{Params (M)} & \hd{GFLOPs} & \hd{Img/s} & \hd{Top-1 (\%)} & \hd{} \\ \gap
MambaVision-T  & 31.8  & 4.4  & 6\,298 & 82.3 & \\
MambaVision-S  & 50.1  & 7.5  & 4\,700 & 83.3 & \\
MambaVision-B  & 97.7  & 15.0 & 3\,670 & 84.2 & \\
MambaVision-L  & 227.9 & 34.9 & 2\,190 & 85.0 & \\ \gap
\hd{COCO backbone} & \hd{Params (M)} & \hd{GFLOPs} & \hd{Box AP} & \hd{Mask AP} & \hd{$\Delta$ box} \\ \gap
ConvNeXt-T     & 86  & 741 & 50.4 & 43.7 & \\
Swin-T         & 86  & 745 & 50.4 & 43.7 & \\
MambaVision-T  & 86  & 740 & 51.1 & 44.3 & $+0.7$ \\ \sgap
ConvNeXt-S     & 108 & 827 & 51.9 & 45.0 & \\
MambaVision-S  & 108 & 828 & 52.3 & 45.2 & $+0.4$ \\ \sgap
ConvNeXt-B     & 146 & 964 & 52.7 & 45.6 & \\
MambaVision-B  & 145 & 964 & 52.8 & 45.7 & $+0.1$ \\
\end{tabularx}
\end{table}

\paragraph{Limitations for this project}
\begin{itemize}
  \item The detection gains over ConvNeXt are small: $+0.7$, $+0.4$ and
        $+0.1$ box AP at the three sizes. They come from single runs with no
        variance reported.
  \item No size-specific metric is given, so its effect on small objects is
        unknown.
  \item Its state-space and attention blocks work only at strides 16 and 32.
        A 16-pixel object barely registers at those strides, so the part of
        the network that should supply global context never sees small
        objects at useful resolution. This is the central mismatch with the
        project.
  \item Throughput is reported only for classification at $224\times224$ on
        an A100. No latency is given for a real-time detector, a consumer GPU
        or batch size one.
\end{itemize}

\paragraph{What to take from it}
Three lessons carry over. Attention or global mixing belongs late in a
stage. A non-SSM branch beside the scan is cheap and helps. The design that
would actually serve small objects, cheap global mixing at strides 4 and 8,
was not attempted.

\section{HPS-DETR: lightweight small-object detection with a transformer}
\label{sec:src-hps}

\paragraph{What it proposes}
HPS-DETR, by Wang and Chen, modifies the real-time transformer detector
RT-DETR-r18 for small objects in aerial images~\citep{wang2025hpsdetr,zhao2024rtdetr}.
It is a method paper, not a survey. It makes four changes.
\begin{itemize}
  \item A lighter backbone block combines partial convolution from
        FasterNet~\citep{chen2023fasternet} with a gated linear unit.
  \item Cascaded group attention, from EfficientViT~\citep{liu2023efficientvit},
        replaces standard self-attention in the encoder. It is applied to
        $P_5$ only.
  \item A new fusion neck brings the stride-4 level into scale-sequence and
        triple-feature fusion. It borrows from ASF-YOLO~\citep{kang2024asfyolo}
        and adds a learned upsampler derived from DySample~\citep{liu2023dysample}.
  \item The box loss combines Inner-IoU and MPDIoU~\citep{zhang2023inneriou,ma2023mpdiou}.
\end{itemize}
The parameter count falls from 19.9M to 15.5M.

\paragraph{Evidence}
Table~\ref{tab:hps} gives the headline results. The authors repeated every
experiment four times and report averages. That is better practice than the
other three papers follow.

\begin{table}[htbp]
\centering
\caption[HPS-DETR: baseline against full model on three datasets]{HPS-DETR
results from the paper, baseline RT-DETR-r18 against the full model. FPS
was measured on an RTX 4090. No model uses pretrained weights.}
\label{tab:hps}
\small
\begin{tabularx}{\textwidth}{@{}LR{1.9cm}R{1.9cm}R{1.9cm}R{1.9cm}@{}}
\hd{Dataset and model} & \hd{AP\textsubscript{50}} & \hd{AP\textsubscript{50:95}} & \hd{AP\textsubscript{S}} & \hd{FPS} \\ \gap
SIMD, RT-DETR-r18 & 77.4 & 62.0 & 25.4 & 70.3 \\
SIMD, HPS-DETR    & 79.8 & 63.5 & 27.9 & 80.7 \\ \sgap
DOTA tiles, RT-DETR-r18 & 85.1 & 60.7 & 37.0 & 68.8 \\
DOTA tiles, HPS-DETR    & 86.8 & 62.5 & 39.2 & 80.2 \\ \sgap
NWPU VHR-10, RT-DETR-r18 & 88.4 & 57.5 & 26.5 & 74.3 \\
NWPU VHR-10, HPS-DETR    & 91.5 & 60.4 & 29.8 & 82.8 \\ \gap
\multicolumn{5}{@{}l}{\textit{Strongest competitors on SIMD}} \\
YOLO11m & 78.8 & 62.4 & & \\
DEIM    & 79.1 & 63.0 & & \\
RT-DETR-r50 & 78.3 & 62.7 & & \\
\end{tabularx}
\end{table}

Its attention ablation is a template worth copying. Holding the rest of the
network fixed, it swaps six attention modules on DOTA. Standard
self-attention gives 84.8 AP\textsubscript{50} at 74.5 FPS. Deformable
attention gives 85.5 at 67.6. HiLo gives 86.0 at 76.6. Cascaded group
attention gives 86.8 at 80.2. A state-space block can be added to such a
table as one more row, so that it competes on equal terms.

\paragraph{Limitations}
\begin{itemize}
  \item Its own loss table suggests that run-to-run variation is about 1 to
        1.5 points. Inner-MPDIoU with ratio 1.0 is mathematically equivalent
        to MPDIoU, yet it scores 1.6 AP\textsubscript{50} lower on SIMD
        (78.2 against 79.8). Several of the claimed gains are no larger than
        this.
  \item The small-object score AP\textsubscript{S} on NWPU swings from 23.2
        to 29.8 with a single constant in the upsampler.
  \item The DOTA experiments use the authors' own tiles and an unstated box
        protocol, so they cannot be compared with the official leaderboard.
  \item Several citations do not point to the works the text attributes to
        them.
  \item The work uses visible light only.
\end{itemize}

\paragraph{What to take from it}
The stride-4 level contributes most of the small-object gain. Global
attention is affordable only at $P_5$, which is precisely where a linear-time
state-space block could be tried at $P_3$ and $P_4$ instead. Repeated runs
are both possible and persuasive.

\section{UAVDet: a CNN--Mamba hybrid for UAV imagery}
\label{sec:src-uavdet}

\paragraph{What it proposes}
UAVDet, by Yang, Guo and Niu, is built on RTMDet-m in MMDetection, not on a
YOLO code base~\citep{yang2026uavdet,lyu2022rtmdet}. It makes three changes.
\begin{itemize}
  \item The backbone is made heavier in its early stages and lighter in its
        late stages. Channels become 192--192--96--96 and block counts
        4--2--2--1.
  \item A cross-stage partial Mamba block (CSPMB) is placed in the neck at
        $P_3$ and $P_4$. It runs a VMamba-style four-direction scan beside a
        convolutional branch.
  \item The neck is redesigned to add a stride-4 head and remove the
        stride-32 head.
\end{itemize}
The paper links to no code.

\paragraph{Evidence}
Table~\ref{tab:uavdet} shows the comparison on the VisDrone test set and the
ablation that explains it. The ablation is the most informative result in
the four papers. On top of the other two changes, the Mamba block adds
0.7~AP while frame rate falls from 76 to 53. A transformer block in the same
position is worse on both counts: 25.6~AP, 46~FPS and seven times the memory.
The authors offer a Mamba-free variant for this reason. They also show that
a plain path-aggregation neck with an added stride-4 head reaches 25.2~AP,
the highest value in their neck study.

\begin{table}[htbp]
\centering
\caption[UAVDet: VisDrone comparison and module ablation]{UAVDet results from
the paper on the VisDrone2019 test set at $640\times640$, measured on an RTX
4090. The upper block compares detectors and the lower block isolates each
module. OB is the optimised backbone. TFPN is the redesigned neck with a
stride-4 head.}
\label{tab:uavdet}
\small
\begin{tabularx}{\textwidth}{@{}LR{1.15cm}R{1.15cm}R{1.15cm}R{1.45cm}R{1.35cm}R{1.0cm}@{}}
\hd{Model} & \hd{AP} & \hd{AP\textsubscript{50}} & \hd{AP\textsubscript{S}} & \hd{Params (M)} & \hd{GFLOPs} & \hd{FPS} \\ \gap
YOLOv8m        & 19.3 & 33.9 & 9.4  & 25.9  & 79.1  & 106 \\
YOLO11m        & 20.6 & 35.7 & 10.7 & 20.1  & 68.2  & 133 \\
Mamba YOLO-M   & 18.3 & 32.2 & 9.0  & 21.8  & 49.6  & 141 \\
RemDet-m       & 23.3 & 39.2 & 11.6 & 23.3  & 34.4  & 188 \\
RT-DETR        & 24.5 & 42.3 & 13.5 & 42    & 136   & 68 \\
RTMDet-m (baseline) & 22.3 & 37.5 & 10.3 & 24.71 & 39.27 & 59 \\
UAVDet         & 26.8 & 44.8 & 15.3 & 3.74\unv & 46.63 & 53 \\ \gap
\multicolumn{7}{@{}l}{\textit{Ablation, each row relative to the baseline}} \\
$+$ OB                  & 25.6 & 43.5 & 14.3 & 23.96 & 42.68 & \\
$+$ CSPMB only          & 22.9 & 38.3 & 10.8 & 29.23 & 43.80 & \\
$+$ TFPN only           & 24.7 & 42.0 & 13.5 & 16.54 & 39.03 & \\
$+$ OB, TFPN            & 26.1 & 44.0 & 15.2 & 3.40\unv & 45.18 & 76 \\
$+$ OB, TFPN, transformer & 25.6 & 42.9 &    & 3.54  & 54.20 & 46 \\
$+$ OB, TFPN, CSPMB     & 26.8 & 44.8 & 15.3 & 3.74\unv & 46.63 & 53 \\
\end{tabularx}
\end{table}

\paragraph{Limitations}
\begin{itemize}
  \item The parameter counts are implausible. The full model is listed at
        3.74M parameters, although its separate components are listed at
        16.54M and 23.96M. The paper does not explain this, and a
        reimplementation should check it.
  \item AP\textsubscript{S} for the baseline is 10.3 in three tables and
        14.5 in two others.
  \item Frame rate is measured with an unstated protocol. The final model
        is slower than its baseline and slower than every YOLO in the
        comparison.
  \item All results come from single runs. A confidence threshold of 0.3 was
        used at inference, which can depress AP if the same threshold was
        applied during evaluation.
  \item The DroneVehicle experiments use the visible images only.
\end{itemize}

\paragraph{What to take from it}
The most important lesson in the four papers comes from this ablation. On
still images, a stride-4 head and early capacity buy several points. The
state-space block buys less than one point and costs about a third of the
frame rate.

\section{CFMW: cross-modality fusion Mamba under adverse weather}
\label{sec:src-cfmw}

\paragraph{What it proposes}
CFMW, by Li and colleagues, targets visible--thermal detection in rain, fog
and snow~\citep{li2025cfmw}. It has two parts. A conditional diffusion model
first restores the weather-degraded visible image. It follows WeatherDiff and
uses DDIM sampling~\citep{ozdenizci2023weatherdiff,ho2020ddpm,song2021ddim}.
A dual-stream YOLOv5-L then extracts features from both images and fuses
them at three backbone stages with a cross-modality fusion Mamba block. That
block does four things in turn.
\begin{itemize}
  \item It swaps half of the channels between the two streams.
  \item It normalises and projects each stream.
  \item It applies VMamba's four-direction scan to each stream.
  \item It gates the outputs and adds them back to both streams through
        crossed residuals.
\end{itemize}
The paper also releases SWVI, a benchmark of 64\,281 paired images with
synthetic rain, fog and snow, built from LLVIP, M3FD, MSRS and FLIR.

\paragraph{Evidence}
Table~\ref{tab:cfmw} gives the main comparison and the efficiency study. The
efficiency figures are the most revealing. The state-space block triples
computation, from 95.6 to 308.6 GFLOPs, for 3.2 points of mAP. The full
model runs at 13.5 frames per second, which is not real-time. On LLVIP, the
thermal-only YOLOv10 scores 62.5 mAP. The plain two-stream YOLOv10 scores
63.4. Fusion therefore adds less than one point on that dataset before any
new module is introduced.

\begin{table}[htbp]
\centering
\caption[CFMW: LLVIP and SWVI comparison and efficiency]{CFMW results read from
Tables IV, V and VII of the paper. mAP is averaged over IoU 0.5 to 0.95. The
two values marked with a dagger disagree with other parts of the paper
(Table~\ref{tab:inconsistencies}).}
\label{tab:cfmw}
\small
\begin{tabularx}{\textwidth}{@{}LR{1.5cm}R{1.5cm}R{1.5cm}R{1.6cm}R{1.4cm}@{}}
\hd{Model, input} & \hd{mAP\textsubscript{50}} & \hd{mAP\textsubscript{75}} & \hd{mAP} & \hd{GFLOPs} & \hd{FPS} \\ \gap
\multicolumn{6}{@{}l}{\textit{LLVIP}} \\
YOLOv10, thermal only & 95.3 & 72.4 & 62.5 & & \\
YOLOv10, two-stream   & 96.1 & 72.7 & 63.4 & & \\
CFT, two-stream       & 97.5 & 72.9 & 63.6 & & \\
CFMW                  & 98.8 & 77.2 & 69.8\unv & & \\ \gap
\multicolumn{6}{@{}l}{\textit{SWVI}} \\
YOLOv10, thermal only & 94.7 & 70.9 & 62.7 & & \\
YOLOv10, two-stream   & 92.1 & 68.2 & 59.3 & & \\
CFT, two-stream       & 94.4 & 69.7 & 59.4 & & \\
CFMW                  & 97.2 & 75.9 & 68.4\unv & & \\ \gap
\multicolumn{6}{@{}l}{\textit{SWVI efficiency study at $640\times640$}} \\
CFT                         & 93.4 & 71.8 & 59.7 & 290.73 & 5.14 \\
CFMW without channel swap and scan & 92.2 & 70.6 & 58.4 & 95.58 & 23.64 \\
CFMW without the scan       & 94.8 & 72.3 & 60.2 & 95.58 & 19.54 \\
CFMW without channel swap   & 96.4 & 75.8 & 62.5 & 308.62 & 15.73 \\
CFMW, full                  & 97.2 & 76.9 & 63.4 & 308.62 & 13.52 \\
\end{tabularx}
\end{table}

\paragraph{Limitations}
\begin{itemize}
  \item The weather is synthetic, and only the visible image is degraded.
        Real fog and rain also affect thermal sensors.
  \item The method depends on pixel-aligned pairs. Channel swapping and
        element-wise addition assume that the two images correspond pixel
        for pixel.
  \item Training the restoration model takes about three days on a 48~GB
        A6000.
  \item Epochs and input size for the detector are not stated in the text.
  \item The authors note that small objects remain a failure case. Objects
        below 2\,500 square pixels make up 26.2\% of SWVI.
  \item Fusion-Mamba, ICAFusion and other contemporaries are not compared.
\end{itemize}

\paragraph{What to take from it}
CFMW is already the journal version of the static design implied by the
project title. It shows how the design is built and evaluated. It also shows
the price: two backbones, a scan at three stages and a frame rate far below
real time. Any new contribution must either avoid that price or buy
something with it that still images cannot.

\section{What the four papers show together}
\label{sec:src-synthesis}

\subsection{Where the gains come from}

Figure~\ref{fig:gains} collects every isolated measurement in the four
papers, and in two closely related works, of what the state-space component
contributes. These are set beside the contributions of other components in
the same papers. The picture is consistent. On still images, the
state-space part adds between $-1.3$ and $+3.2$ points, usually below one
point and often at a large cost in speed or computation. The stride-4 head
and redistributed early capacity add more, and they do so cheaply.

\begin{figure}[htbp]
\centering
\begin{tikzpicture}
  \begin{axis}[
      width=8.2cm, height=12.4cm,
      xbar, bar width=7pt,
      xmin=-2, xmax=4.2,
      ymin=0.4, ymax=12.6,
      axis lines=left, axis line style={draw=cNeutral, line width=0.5pt},
      y axis line style={draw=none},
      ytick={1,...,12}, yticklabels={},
      ytick style={draw=none},
      xtick={-1,0,1,2,3,4},
      tick label style={font=\scriptsize, text=muted},
      xlabel={change in AP or mAP, points},
      xlabel style={font=\footnotesize, text=muted},
      clip=false,
      enlarge y limits=0,
      extra x ticks={0}, extra x tick style={grid=major, major grid style={draw=faint, line width=0.4pt}, tick label style={opacity=0}},
    ]
    % state-space components (purple), rows 11 to 5
    \addplot[draw=none, fill=cTime!70] coordinates {(3.2,12) (0.9,11) (0.7,10) (0.7,9) (0.6,8) (0.1,7) (-1.3,6)};
    % other components (grey), rows 4 to 1
    \addplot[draw=none, fill=cNeutral!55] coordinates {(3.5,4) (3.3,3) (2.4,2) (1.5,1)};

    % row labels on the left
    \def\L#1#2{\node[anchor=east, font=\scriptsize, text=ink, align=right] at (axis cs:-2.15,#1) {#2};}
    \L{12}{CFMW: scan in fusion block, SWVI}
    \L{11}{MambaPSA in YOLO26, best placement, VOC}
    \L{10}{UAVDet: Mamba added to CSP neck, VisDrone}
    \L{9}{MambaVision-T against ConvNeXt-T, COCO}
    \L{8}{UAVDet: CSP-Mamba block alone, VisDrone}
    \L{7}{MambaVision-B against ConvNeXt-B, COCO}
    \L{6}{Mamba YOLO-L against YOLO11-l, COCO}
    \L{4}{CFMW: weather restoration alone, SWVI}
    \L{3}{UAVDet: early-heavy backbone, VisDrone}
    \L{2}{UAVDet: stride-4 neck and head, VisDrone}
    \L{1}{HPS-DETR: four changes together, SIMD}

    % cost annotations on the right
    \def\C#1#2#3{\node[anchor=west, font=\scriptsize\itshape, text=muted] at (axis cs:#1,#2) {#3};}
    \C{3.35}{12}{GFLOPs $\times3.2$}
    \C{1.05}{11}{range $-0.1$ to $+0.9$}
    \C{0.85}{10}{76 to 53 FPS}
    \C{0.85}{9}{single run}
    \C{0.75}{8}{$+4.5$M params}
    \C{0.25}{7}{single run}
    \C{0.15}{6}{similar FLOPs}
    \C{3.65}{4}{3 days to train}
    \C{3.45}{3}{GFLOPs $+3.4$}
    \C{2.55}{2}{GFLOPs $-0.2$}
    \C{1.65}{1}{FPS $+10$}

    % group headings
    \node[anchor=west, font=\footnotesize\scshape, text=cTime] at (axis cs:-1.95,12.75) {state-space component};
    \node[anchor=west, font=\footnotesize\scshape, text=cNeutral] at (axis cs:-1.95,4.85) {other components};
  \end{axis}
\end{tikzpicture}

\vspace{2mm}
\caption[Measured contribution of the state-space component against other
components]{Measured contribution of the state-space component (purple)
against other components (grey), from the ablations of the source papers and
two related works~\citep{li2025cfmw,yang2026uavdet,hatamizadeh2025mambavision,wang2025hpsdetr,mambapsa2026,wang2025mambayolo,yolo11_2024}.
Each bar isolates one change on one dataset. Notes on the right give the
cost reported with it. On still images, the stride-4 pathway and the
redistribution of early capacity contribute more than the state-space block,
and at lower cost.}
\label{fig:gains}
\end{figure}


\subsection{What none of them do}

Table~\ref{tab:capabilities} lists the capabilities a journal reviewer would
look for. None of the four papers processes video, carries memory across
frames, tests robustness to sensor failure or measures latency on an edge
device. Only one repeats its experiments. These gaps define where a new
contribution can stand.

\begin{table}[htbp]
\centering
\caption[Capabilities covered by the four source papers]{Capabilities covered
by the four source papers. Entries are as reported in each paper.}
\label{tab:capabilities}
\small
\begin{tabularx}{\textwidth}{@{}LC{2.0cm}C{2.0cm}C{2.0cm}C{2.0cm}@{}}
\hd{Capability} & \hd{MambaVision} & \hd{HPS-DETR} & \hd{UAVDet} & \hd{CFMW} \\ \gap
state-space component      & yes & no  & yes, neck & yes, fusion \\
thermal input              & no  & no  & no  & yes \\
small-object metric        & no  & AP\textsubscript{S} & AP\textsubscript{S}, AP\textsubscript{T} & size analysis \\
video or temporal memory   & no  & no  & no  & no \\
robustness to sensor faults & no & no  & no  & weather only \\
repeated runs              & no  & four & no  & no \\
latency at batch one       & no  & yes, 4090 & yes, 4090 & yes, GPU unstated \\
edge deployment            & no  & no  & no  & no \\
code released              & yes & no  & no  & announced \\
\end{tabularx}
\end{table}

\subsection{Reporting problems}

All four papers contain internal inconsistencies (Table~\ref{tab:inconsistencies}).
None of them changes the main conclusions. Together they argue for two
habits. Every baseline in the new work should be re-run under one protocol,
not copied from papers. Every table should be generated directly from logged
runs (Section~\ref{sec:nb-results}).

\begin{table}[htbp]
\centering
\caption[Inconsistencies found in the source papers]{Inconsistencies found in
the source papers. Values are quoted as printed.}
\label{tab:inconsistencies}
\small
\begin{tabularx}{\textwidth}{@{}P{2.3cm}L@{}}
\hd{Paper} & \hd{Inconsistency} \\ \gap
MambaVision & Ablation text gives a baseline of 80.9\% top-1 and 40.2 mask AP; the table gives 80.5\% and 40.4. \\
            & Claims 56\% fewer GFLOPs than MaxViT-B; 15.0 against 23.4 is 36\% fewer. \\
            & The output projection in Eq.~7 has the wrong input dimension. \\ \sgap
HPS-DETR    & Inner-MPDIoU at ratio 1.0, equivalent to MPDIoU, differs from it by 1.6 AP\textsubscript{50}. \\
            & The new backbone block is said to speed the model up, yet DOTA FPS falls from 68.8 to 65.2. \\
            & References [4], [20], [30] and [34] do not match the works described. \\ \sgap
UAVDet      & The full model has 3.74M parameters although its components have 16.54M and 23.96M. \\
            & Baseline AP\textsubscript{S} is 10.3 in Tables 3, 5 and 6 and 14.5 in Tables 7 and 8. \\
            & CSPMB at $L_3$--$L_5$ has 29.23M parameters in Table 5 and 29.41M in Table 6. \\ \sgap
CFMW        & LLVIP mAP is 69.8 in Table IV and 64.8 in the text. \\
            & The full model scores 68.4 mAP on SWVI in Table V and 63.4 in Tables VI and VII. \\
            & The SWVI splits sum to 59\,990 images, not the stated 64\,281. \\
            & A 51.2\% memory saving is claimed; 18.0 against 35.2\,GB is 48.9\%. \\
            & Three fusion blocks are described, but performance peaks at eight in Fig.~7. \\
\end{tabularx}
\end{table}
